This blueprint covers every numbered result of \emph{On attention heads and
bilinear forms}
(\href{https://arxiv.org/abs/2609.22990v1}{arXiv:2609.22990v1}), all of which
are formalized in Lean: the lemmas,
propositions and theorems of the main text and of Appendices A, B and C,
together with the definitions they are stated in. Numbering follows arXiv:2609.22990v1;
Theorems 1 and 2 of the introduction restate Theorems 3 and 5. Definitions
numbered A1--A3 are auxiliary: they collect notation that the paper introduces
in the running text. Each result has a proof outline and links to its Lean
declarations; in the \href{dep_graph_document.html}{dependency graph}, green
means that the statement and its proof are both formalized. The empirical
observations and the two remarks of the main text, which are stated without
proof, are not formalized.

\chapter{Attention functions and transformers}

Throughout, $C$ is a finite-dimensional real vector space. Smooth means
$C^\infty$, and $\Delta_{k-1}\subset\R^k$ is the standard simplex.

\begin{definition}[Attention function]
\label{def:attention-function}
\lean{Appendices.MainText.IsPolyFun, Appendices.MainText.IsAttentionFunction}
\leanok
A smooth map $a\colon C^k\to C$ is an attention function if there are a
smooth $\omega\colon C^k\to\Delta_{k-1}$, a smooth convex
$\phi\colon\R^k\to\R$ and a polynomial map $\sigma\colon C^k\to\R^k$ with
\[
 \omega=(\nabla\phi)\circ\sigma,\qquad a(x_1,\ldots,x_k)=\sum_{j=1}^k\omega_j(x)\,x_j.
\]
The data $\omega,\phi,\sigma$ are carried along with $a$. A polynomial map is
one whose components are polynomials in the coordinates of $x_1,\ldots,x_k$.
\end{definition}

\begin{definition}[Biaffine and bilinear attention functions]
\label{def:bilinear}
\lean{Appendices.MainText.IsBilinearScore}
\leanok
\uses{def:attention-function}
An attention function whose scoring function has the form
$\sigma(x)=(P_{k-1}(x_1,x_k),\ldots,P_0(x_k,x_k))$ for polynomials
$P_d\colon C^2\to\R$ of bidegree at most $(1,1)$ is biaffine; it is bilinear
if every $P_d$ is a bilinear form.
\end{definition}

\begin{lemma}[Biaffine decomposition]
\label{lem:biaffine-decomposition}
\lean{Appendices.MainText.IsPolyBidegreeLeOneOne, Appendices.MainText.biaffine_components, Appendices.MainText.biaffine_decomposition}
\leanok
Every polynomial $P\colon C\times C\to\R$ of bidegree at most $(1,1)$ has a
unique decomposition $P(u,v)=B(u,v)+L(u)+R(v)+D$ with $B$ bilinear, $L,R$
linear and $D\in\R$, namely
\[
 D=P(0,0),\quad L(u)=P(u,0)-P(0,0),\quad R(v)=P(0,v)-P(0,0),
\]
\[
 B(u,v)=P(u,v)-P(u,0)-P(0,v)+P(0,0).
\]
\end{lemma}
\begin{proof}
\leanok
The paper calls this elementary and gives no proof. Each monomial of
bidegree at most $(1,1)$ is a product of an affine function of $u$ and an
affine function of $v$; expanding the product shows that the four
expressions are bilinear, linear, linear and constant, and they add up to
$P$. Evaluating any decomposition at $(0,0)$, $(u,0)$ and $(0,v)$ gives the
formulas, hence uniqueness.
\end{proof}

\begin{definition}[RoPE-like]
\label{def:rope-like}
\lean{Appendices.MainText.ropeL, Appendices.MainText.ropeL0Equiv, Appendices.MainText.ropeT, Appendices.MainText.IsRoPELike}
\leanok
\uses{def:bilinear}
Let a bilinear attention function on a head space $H$ have forms
$B_{k-1},\ldots,B_0$, let $L_j\colon H\to H^\vee$, $L_j(v)(u)=B_j(u,v)$,
assume $B_0$ nondegenerate, and set $T_d=L_0^{-1}L_d$. The function is
RoPE-like if $T_{d+e}=T_dT_e$ whenever $d,e\ge0$ and $d+e<k$.
\end{definition}

\begin{proposition}[Classification of RoPE-like attention functions]
\label{prop:rope-like}
\lean{Appendices.MainText.isRoPELike_iff, Appendices.MainText.eq_ropeT_one_of_forall}
\leanok
\uses{def:rope-like}
With $B_0$ nondegenerate, the function is RoPE-like if and only if there is
$S\in\operatorname{End}(H)$ with $T_d=S^d$ for all $0\le d<k$. When $S$
exists it is $T_1=L_0^{-1}L_1$.
\end{proposition}
\begin{proof}
\leanok
\uses{def:rope-like}
The RoPE-like condition gives $L_{d+e}=L_dL_0^{-1}L_e$, in particular
$L_{d+1}=L_dS$ with $S=T_1$. By induction $L_d=L_0S^d$, that is
$T_d=S^d$. Conversely $S^{d+e}=S^dS^e$.
\end{proof}

\begin{definition}[Attention head]
\label{def:head}
\lean{Appendices.MainText.head}
\leanok
\uses{def:attention-function}
For attention functions $a_k\colon C^k\to C$, $k=2,\ldots,\ell$, the causal
attention head is
$A(x_1,\ldots,x_\ell)=(x_1,a_2(x_1,x_2),\ldots,a_\ell(x_1,\ldots,x_\ell))$.
\end{definition}

\begin{definition}[Transformer]
\label{def:transformer}
\lean{Appendices.MainText.transformer}
\leanok
\uses{def:head}
For a head $A$, a linear value map $V\colon C\to C$ and a map
$M\colon C\to C$, the single-head transformer is $T=T_M\circ T_A$ with
$T_A(x)_i=x_i+V(A(x)_i)$ and $T_M(y)_i=y_i+M(y_i)$.
\end{definition}

\begin{proposition}[Transformers are near-identity]
\label{prop:near-identity}
\lean{Appendices.MainText.IsAttentionOutput, Appendices.MainText.attentionBlock, Appendices.MainText.norm_attentionBlock_le, Appendices.MainText.near_identity}
\leanok
\uses{def:transformer}
Fix a norm on $C$ and put $\norm{x}_\infty=\max_k\norm{x_k}$. If
$\norm{M(y)}\le L\norm y$ for all $y$ and $V=\sum_i\norm{V_i}$ (operator
norms), then $\norm{T(x)-x}_\infty\le(L+V+LV)\norm{x}_\infty$.
\end{proposition}
\begin{proof}
\leanok
\uses{def:transformer}
Each component of a head is a convex combination of its inputs, so its norm
is at most $\norm x_\infty$; the attention block is then bounded by
$V\norm x_\infty$. The $k$-th component of $T(x)-x$ is
$A(x)_k+M(x_k+A(x)_k)$, whose norm is at most
$V\norm x_\infty+L(1+V)\norm x_\infty$. The Lean statement assumes only that
each head returns convex combinations of its inputs, which is what the proof
uses.
\end{proof}

\begin{proposition}[Bilinear attention recovers QKV attention]
\label{prop:homogenization}
\lean{Appendices.MainText.qkvScore, Appendices.MainText.homForm, Appendices.MainText.homForm_apply, Appendices.MainText.homScore, Appendices.MainText.dehom, Appendices.MainText.homScore_eq, Appendices.MainText.isPolyFun_bilinear, Appendices.MainText.qkv_homogenizes_to_bilinear}
\leanok
\uses{def:bilinear,def:transformer}
Let $\widehat C=C\oplus\R e$ and $\iota(x)=x+e$. For every single-head QKV
transformer $T\colon C^\ell\to C^\ell$, with scores
$\langle Q(x_k+p_k),K(x_j+p_j)\rangle/\sqrt n$, there is a single-head
bilinear transformer $\widehat T$ on $\widehat C$ with the same potentials
such that $\widehat T(\iota x_1,\ldots,\iota x_\ell)=\iota^\ell(T(x_1,\ldots,x_\ell))$.
\end{proposition}
\begin{proof}
\leanok
\uses{def:bilinear,def:transformer}
Define
$\widehat B(u+\alpha e,v+\beta e)=\langle Q(v+\beta p_k),K(u+\alpha p_j)\rangle/\sqrt n$,
which expands to the paper's four-term formula. It agrees with the QKV score
on the chart $C+e$, so the weights agree there and
$\sum_j\omega_j(x_j+e)=a_k(x)+e$. Extend $V$ and $M$ by ignoring the new
coordinate. The homogenized weights are simplex-valued on all of
$\widehat C^k$ because $\widehat\sigma(u+\alpha e,\ldots)=\sigma(x)$ with
$x_j=u_j+(\alpha_j-1)p_j$; this also gives their smoothness. The Lean
statement allows any potentials, softmax included.
\end{proof}

\chapter{Pairing and balance}

\setcounter{definition}{0}
\renewcommand{\thedefinition}{A\arabic{definition}}
\begin{definition}[Symmetric and antisymmetric parts]
\label{def:parts}
\lean{Appendices.frobSq, Appendices.symPart, Appendices.skewPart}
\leanok
For $L\in M_N(\R)$, put
\[
 S=\tfrac12(L+L^T),\qquad T=\tfrac12(L-L^T),\qquad
 \norm{L}^2=\tr(L^TL).
\]
Then $L=S+T$, with $S$ symmetric and $T$ antisymmetric.
\end{definition}
\renewcommand{\thedefinition}{\arabic{definition}}

\begin{lemma}[Pairing bound]
\label{lem:pairing-bound}
\lean{Appendices.MainText.lamPlusVec, Appendices.MainText.lamMinusVec, Appendices.MainText.pairingScoreSym, Appendices.MainText.frobSq_eq_norm_sq_add_norm_sq, Appendices.MainText.pairing_bound}
\leanok
\uses{def:parts}
Let $S$ be symmetric, let $\lambda_\pm$ be its sorted positive eigenvalues and
absolute values of negative eigenvalues, padded with zeros to equal length,
and let $\mathcal P=1-\norm{\lambda_+-\lambda_-}/\norm S$. If
$\lambda_+\ne0$ and $\rho=\norm{\lambda_-}/\norm{\lambda_+}$, then
\[
 \mathcal P=1-\sqrt{1-\frac{2\langle\lambda_+,\lambda_-\rangle}{\norm S^2}}
 \le 1-\frac{|1-\rho|}{\sqrt{1+\rho^2}},
\]
with equality if and only if $\lambda_-=\rho\lambda_+$.
\end{lemma}
\begin{proof}
\leanok
\uses{def:parts}
Since $\norm S^2=\norm{\lambda_+}^2+\norm{\lambda_-}^2$, expanding
$\norm{\lambda_+-\lambda_-}^2$ gives the equality. The triangle inequality
gives $\norm{\lambda_+-\lambda_-}\ge|1-\rho|\norm{\lambda_+}$, with equality
if and only if one list is a nonnegative multiple of the other; divide by
$\norm S=\sqrt{1+\rho^2}\norm{\lambda_+}$.
\end{proof}

\begin{lemma}[Rank balance]
\label{lem:rank-balance}
\lean{Appendices.MainText.finrank_le_of_posDefOn, Appendices.MainText.rank_balance}
\leanok
\uses{def:parts}
If $L\in M_N(\R)$ has rank at most $n$, the symmetric part of $L$ has at
most $n$ positive and at most $n$ negative eigenvalues, counted with
multiplicity.
\end{lemma}
\begin{proof}
\leanok
\uses{def:parts}
The quadratic form of $S$ vanishes on $\ker L$, of dimension at least
$N-n$. The span of the eigenvectors with positive eigenvalues meets
$\ker L$ only in $0$, so its dimension $n_+$ is at most $n$. For $n_-$,
apply the same argument to $-L$.
\end{proof}

\setcounter{theorem}{2}
\begin{theorem}[Heads have balanced attention]
\label{thm:balanced}
\lean{Appendices.MainText.stacked, Appendices.MainText.Nondegenerate, Appendices.MainText.headForm, Appendices.MainText.radical_eq_ker, Appendices.MainText.finrank_radical, Appendices.MainText.gram, Appendices.MainText.signature, Appendices.MainText.gram_eq_symPart, Appendices.MainText.ker_gram_eq_map, Appendices.MainText.heads_have_balanced_attention}
\leanok
\uses{lem:rank-balance}
Let $Q,K\colon C\to H$ with $\dim H=n$ and $\dim C=N$. If the stacked map
$J=(Q,K)\colon C\to H\oplus H$ is surjective, the symmetric form
$S(x,y)=\tfrac12(\langle Kx,Qy\rangle+\langle Qx,Ky\rangle)$ has signature
$(n_+,n_-,n_0)=(n,n,N-2n)$. In Lean the signature is read off the Gram
matrix of $S$ in an arbitrary basis of $C$.
\end{theorem}
\begin{proof}
\leanok
\uses{lem:rank-balance}
The Gram matrix of $2S$ is the symmetric part of $2K^TQ$, which has rank at
most $n$, so Lemma 3 gives $n_\pm\le n$. The radical of $S$ is $\ker J$:
if $S(x,\cdot)=0$ then $B(Jx,\cdot)$ vanishes on $J(C)=H\oplus H$ for the
nondegenerate form $B((a,b),(a',b'))=\langle a,b'\rangle+\langle b,a'\rangle$.
By surjectivity $\dim\ker J=N-2n$, so $S$ has rank $2n=n_++n_-$, which
forces $n_+=n_-=n$.
\end{proof}

\chapter{Profiles of bilinear forms}

All matrices are real. Matrix norms are Frobenius norms. The Lean
development uses their squares to avoid unnecessary square roots.

\setcounter{definition}{5}
\begin{definition}[Profile of a bilinear form]
\label{def:profile}
\lean{Appendices.sortDesc, Appendices.alphaList, Appendices.betaList, Appendices.profileA, Appendices.profileB, Appendices.profileC, Appendices.profileD}
\leanok
\uses{def:parts}
For $L\ne0$, let $\alpha$ and $\beta$ be the decreasing lists of positive
eigenvalues of $S$ and absolute values of its negative eigenvalues,
padded with zeros to equal length. Define $\pi(L)=(a,b,c,d)$ by
\[
 a=\frac{\norm{T}^2}{\norm{L}^2},\qquad
 b=\frac{2\langle\alpha,\beta\rangle}{\norm{L}^2},\qquad
 c=\frac{\norm{(\alpha-\beta)_+}^2}{\norm{L}^2},\qquad
 d=\frac{\norm{(\alpha-\beta)_-}^2}{\norm{L}^2}.
\]
Positive and negative parts are taken entrywise. Lean pads both lists
to length $N$ and includes the denominator in each definition.
\end{definition}

\begin{lemma}[Profile scale]
\label{lem:profile-scale}
\lean{Appendices.MainText.profileA_symPart, Appendices.MainText.one_sub_profileA, Appendices.MainText.profile_scale}
\leanok
\uses{def:profile,lem:energy}
If the symmetric part $S$ of $L$ is nonzero, $\pi(L)=(a,b,c,d)$ and
$\pi(S)=(0,b_S,c_S,d_S)$, then $(b,c,d)=(1-a)(b_S,c_S,d_S)$.
\end{lemma}
\begin{proof}
\leanok
\uses{def:profile,lem:energy}
Both triples are the same expressions in the eigenvalues of $S$, divided by
$\norm L^2$ and by $\norm S^2$ respectively, and
$\norm S^2/\norm L^2=1-\norm T^2/\norm L^2=1-a$.
\end{proof}

\setcounter{proposition}{3}
\begin{proposition}[Profiles of rank-one matrices]
\label{prop:rank-one}
\lean{Appendices.MainText.profile, Appendices.MainText.Theta, Appendices.MainText.profile_conj, Appendices.MainText.exists_orthogonal_cols, Appendices.MainText.profile_smul_vecMulVec, Appendices.MainText.exists_smul_vecMulVec_of_rank_eq_one, Appendices.MainText.profile_rank_one, Appendices.MainText.profiles_rank_one_eq_Theta}
\leanok
\uses{def:profile,lem:energy}
If $L$ has rank one and unit norm, then with $t=\tr(L)$,
\[
 \pi(L)=\Bigl(\frac{1-t^2}2,\frac{1-t^2}2,(t_+)^2,(t_-)^2\Bigr).
\]
For $N\ge2$ the set of profiles of rank-one matrices is
$\Theta=\{a=b,\ cd=0\}\subset\Delta_3$. (For $N=1$ only the two vertices
$c=1$ and $d=1$ occur, so the set statement needs $N\ge2$.)
\end{proposition}
\begin{proof}
\leanok
\uses{def:profile,lem:energy}
Write $L=uv^T$ with unit vectors, so $t=\langle u,v\rangle$; the profile is
invariant under orthogonal changes of coordinates. If $u=\pm v$, bring $u$
to the first basis vector. Otherwise $e=(u+v)/\sqrt{2+2t}$ and
$f=(v-u)/\sqrt{2-2t}$ are orthonormal; in a basis extending them $L$ is the
block $\tfrac12\bigl(\begin{smallmatrix}1+t&\sqrt{1-t^2}\\-\sqrt{1-t^2}&t-1\end{smallmatrix}\bigr)$,
whose symmetric part is $\operatorname{diag}((1+t)/2,(t-1)/2)$. Read off
the profile. Every point of $\Theta$ arises from $u=e_1$,
$v=te_1+\sqrt{1-t^2}\,e_2$.
\end{proof}

\setcounter{theorem}{4}
\begin{theorem}[Accumulation]
\label{thm:accumulation}
\lean{Appendices.MainText.profile_of_spectrum, Appendices.MainText.accumulation_estimates, Appendices.MainText.accumulation, Appendices.MainText.limitProfile, Appendices.MainText.projP, Appendices.MainText.limitProfile_mem_projP_Theta, Appendices.MainText.limitB, Appendices.MainText.limitProfile_edge_d, Appendices.MainText.limitProfile_edge_c, Appendices.MainText.limitB_strictMonoOn, Appendices.MainText.limitB_strictAntiOn, Appendices.MainText.limitB_image_Ioc, Appendices.MainText.limitB_image_Ici}
\leanok
\uses{lem:proportional,prop:rank-one}
Let $X>0$ have finite nonzero second moment, fix $u,v>0$ and $\rho=u/v$. Let
$N_n\ge2n$ and let $S_n\in M_{N_n}(\R)$ be symmetric with eigenvalues
$vX_1,\ldots,vX_n$, $-uX'_1,\ldots,-uX'_n$ and zeros, where the $X_i,X'_i$
are independent copies of $X$. Then almost surely
\[
 \pi(S_n)\longrightarrow\frac1{1+\rho^2}\bigl(0,\,2\rho,\,(1-\rho)_+^2,\,(\rho-1)_+^2\bigr)\in p(\Theta),
\]
where $p(a,b,c,d)=(0,a+b,c,d)$. As $\rho$ ranges over $(0,1]$ the limit
traces the edge $\{d=0\}$, and over $[1,\infty)$ the edge $\{c=0\}$.
\end{theorem}
\begin{proof}
\leanok
\uses{lem:proportional,prop:rank-one}
Write $\beta_n=\rho\alpha_n+\varepsilon_n$ for the sorted lobes; by Lemma 6,
$\delta_n=\norm{\varepsilon_n}/\norm{\alpha_n}\to0$ almost surely. By
Cauchy--Schwarz, the triangle inequality and the $1$-Lipschitz property of
$x\mapsto x_\pm$, the ratios $\langle\alpha,\beta\rangle/\norm\alpha^2$,
$\norm\beta^2/\norm\alpha^2$ and $\norm{(\alpha-\beta)_\pm}/\norm\alpha$ lie
within $\delta_n$, $\delta_n^2+2\rho\delta_n$ and $\delta_n$ of $\rho$,
$\rho^2$ and $(1-\rho)_\pm$. Divide by
$\norm{S_n}^2/\norm{\alpha_n}^2=1+\norm{\beta_n}^2/\norm{\alpha_n}^2$. On
each edge, $b(\rho)=2\rho/(1+\rho^2)$ is a strictly monotone bijection onto
$(0,1]$.
\end{proof}

\chapter{Appendix A: Random matrices}

These results give the random baselines for the symmetric energy of an
attention head. The distinction between a ratio of expectations and
the expectation of a ratio matters: the latter uses Gaussian symmetry.

\begin{lemma}[Symmetric and skew energy identities]
\label{lem:energy}
\lean{Appendices.frobSq_eq_frobSq_symPart_add_frobSq_skewPart, Appendices.frobSq_symPart, Appendices.frobSq_skewPart, Appendices.frobSq_symPart_div, Appendices.frobSq_skewPart_div}
\leanok
\uses{def:parts}
For any $L\in M_N(\R)$,
\[
 \norm{L}^2=\norm{S}^2+\norm{T}^2,\qquad
 \norm{S}^2=\tfrac12\norm{L}^2+\tfrac12\tr(L^2),\qquad
 \norm{T}^2=\tfrac12\norm{L}^2-\tfrac12\tr(L^2).
\]
For $L\ne0$, division gives
\[
 \frac{\norm{S}^2}{\norm{L}^2}=\frac12+\frac{\tr(L^2)}{2\norm{L}^2},
 \qquad
 \frac{\norm{T}^2}{\norm{L}^2}=\frac12-\frac{\tr(L^2)}{2\norm{L}^2}.
\]
\end{lemma}
\begin{proof}
\leanok
\uses{def:parts}
Expand the squared Frobenius norms entrywise, using invariance under
transpose and $\sum_{i,j}L_{ij}L_{ji}=\tr(L^2)$. Adding the two
identities gives the orthogonal decomposition of the energy.
\end{proof}

\begin{corollary}[Iid square random matrix baseline]
\label{cor:iid}
\lean{Appendices.integral_frobSq_iid, Appendices.integral_frobSq_symPart_iid, Appendices.integral_frobSq_skewPart_iid, Appendices.integral_frobSq_symPart_div_integral_frobSq_iid, Appendices.integral_frobSq_skewPart_div_integral_frobSq_iid, Appendices.integral_frobSq_symPart_div_frobSq_gaussian, Appendices.integral_frobSq_skewPart_div_frobSq_gaussian}
\leanok
\uses{def:parts}
Let the entries of $L\in M_N(\R)$ be independent, mean zero, and
square-integrable with common variance $\sigma^2$. Then
\[
 \E\norm{L}^2=N^2\sigma^2,\quad
 \E\norm{S}^2=\frac{N(N+1)}2\sigma^2,\quad
 \E\norm{T}^2=\frac{N(N-1)}2\sigma^2.
\]
For $N>0$ and $\sigma\ne0$, the ratios to $\E\norm{L}^2$ are
$(N+1)/(2N)$ and $(N-1)/(2N)$ respectively. If the entries are iid
nondegenerate centered Gaussian, then also
\[
 \E\left[\frac{\norm{S}^2}{\norm{L}^2}\right]=\frac{N+1}{2N},
 \qquad
 \E\left[\frac{\norm{T}^2}{\norm{L}^2}\right]=\frac{N-1}{2N}.
\]
\end{corollary}
\begin{proof}
\leanok
\uses{lem:energy}
Independence and centering kill the off-diagonal cross terms, giving
$\E\tr(L^2)=N\sigma^2$. Apply the energy identities. In the Gaussian
case, symmetry of the joint law makes the expected normalized squared
coordinate equal to $1/N^2$, while sign symmetry cancels the remaining
normalized cross terms.
\end{proof}

\begin{corollary}[Random QK-product baseline]
\label{cor:qk}
\lean{Appendices.integral_frobSq_qk, Appendices.integral_trace_mul_self_qk, Appendices.integral_frobSq_symPart_qk, Appendices.integral_frobSq_skewPart_qk, Appendices.integral_frobSq_symPart_div_integral_frobSq_qk, Appendices.integral_frobSq_skewPart_div_integral_frobSq_qk}
\leanok
\uses{def:parts}
Let $W_K,W_Q\in M_{n\times N}(\R)$ have jointly independent,
mean-zero, square-integrable entries, with variances $\sigma_K^2$ and
$\sigma_Q^2$ respectively. For $L=W_K^TW_Q$,
\[
 \E\norm{L}^2=nN^2\sigma_K^2\sigma_Q^2,\qquad
 \E\norm{S}^2=\frac{nN(N+1)}2\sigma_K^2\sigma_Q^2,\qquad
 \E\norm{T}^2=\frac{nN(N-1)}2\sigma_K^2\sigma_Q^2.
\]
For positive dimensions and nonzero variances, the ratios of the last
two expectations to the first are $(N+1)/(2N)$ and $(N-1)/(2N)$.
\end{corollary}
\begin{proof}
\leanok
\uses{lem:energy}
Expand $L_{ab}=\sum_r(W_K)_{ra}(W_Q)_{rb}$. Independence gives
$\E L_{ab}^2=n\sigma_K^2\sigma_Q^2$ and
$\E\tr(L^2)=nN\sigma_K^2\sigma_Q^2$. Sum the entries and use Lemma 5.
\end{proof}

\begin{proposition}[Expected symmetric energy for random low-rank QK products]
\label{prop:low-rank}
\lean{Appendices.expected_symmetric_energy_low_rank_qk, Appendices.expected_skew_energy_low_rank_qk}
\leanok
\uses{def:parts}
Let $1\le n\le N$ and let $W_K,W_Q\in M_{n\times N}(\R)$ have
jointly independent centered Gaussian entries, iid within each matrix,
with respective nonzero variances $v_K,v_Q$. For $L=W_K^TW_Q$,
\[
 \E\left[\frac{\norm{S}^2}{\norm{L}^2}\right]=\frac{N+1}{2N},
 \qquad
 \E\left[\frac{\norm{T}^2}{\norm{L}^2}\right]=\frac{N-1}{2N},
\]
independently of the head dimension $n$.
\end{proposition}
\begin{proof}
\leanok
\uses{lem:energy}
First take equal variances. Sign changes and permutations of Gaussian
columns preserve the joint law. Sign symmetry cancels the off-diagonal
terms of $\tr(L^2)/\norm{L}^2$; permutation symmetry makes the expected
normalized squared entries equal. Thus its expectation is $1/N$, and
Lemma 5 gives the result. The denominator is nonzero almost surely.
Rescale $W_K$ to reduce unequal variances to the equal-variance case;
the normalized energies are unchanged by this rescaling.
\end{proof}

\begin{lemma}[Proportional lobes]
\label{lem:proportional}
\lean{Appendices.proportional_lobes}
\leanok
\uses{def:profile}
Let $X>0$ almost surely, with finite nonzero second moment, and fix
$u,v>0$, $\rho=u/v$. Let $\alpha_n$ be the decreasing sorting of
$vX_1,\ldots,vX_n$ and $\beta_n$ the decreasing sorting of
$uX'_1,\ldots,uX'_n$, where the two sequences together consist of
independent copies of $X$. Then
\[
 \frac{\norm{\beta_n-\rho\alpha_n}}{\norm{\alpha_n}}
 \longrightarrow 0\qquad\text{almost surely as }n\longrightarrow\infty.
\]
\end{lemma}
\begin{proof}
\leanok
\uses{def:profile}
Reduce to two samples with the same scale. For sorted nonnegative
lists, a layer-cake identity expresses their inner product using the
minimum of two empirical survival counts. The strong law controls
these counts and the empirical second moments. Fatou's lemma bounds
the limiting inner product from below by the common second moment.
Expanding the squared difference gives convergence to zero after
division by $n$; the squared denominator divided by $n$ converges to
the strictly positive second moment. Restore the scales $u,v$.
\end{proof}

\chapter{Appendix B: Balanced bimodal spectra and simplex faces}

\setcounter{definition}{1}
\renewcommand{\thedefinition}{A\arabic{definition}}
\begin{definition}[Pairing geometry]
\label{def:balanced}
\lean{Appendices.MemLambdaSet, Appendices.lamPlus, Appendices.lamMinus, Appendices.imbalance, Appendices.statC, Appendices.pairingScore, Appendices.iotaMap}
\leanok
For $m\ge1$, let $\Lambda_m$ consist of the unit vectors
$\lambda_1\le\cdots\le\lambda_{2m}$ with
$\lambda_m<0<\lambda_{m+1}$. Set
\[
 (\lambda_+)_j=\lambda_{2m+1-j},\qquad
 (\lambda_-)_j=-\lambda_j\quad(1\le j\le m).
\]
Define
\[
 \mathcal I=\norm{\lambda_+}^2-\norm{\lambda_-}^2,\quad
 C=\langle\lambda_+,\lambda_-\rangle,\quad
 \mathcal P=1-\norm{\lambda_+-\lambda_-},\quad
 (\iota\lambda)_i=-\lambda_{2m+1-i}.
\]
\end{definition}

\begin{lemma}[Bimodal identities]
\label{lem:bimodal}
\lean{Appendices.sum_sq_lamPlus_eq, Appendices.sum_sq_lamMinus_eq, Appendices.two_mul_norm_lamPlus_mul_norm_lamMinus, Appendices.sum_sq_lamPlus_sub_lamMinus, Appendices.pairingScore_eq, Appendices.sum_sq_complex, Appendices.statC_nonneg, Appendices.two_mul_statC_le_one, Appendices.pairingScore_nonneg, Appendices.pairingScore_le_one}
\leanok
\uses{def:balanced}
For $\lambda\in\Lambda_m$,
\[
 \norm{\lambda_+}^2=\tfrac12(1+\mathcal I),\qquad
 \norm{\lambda_-}^2=\tfrac12(1-\mathcal I),\qquad
 2\norm{\lambda_+}\norm{\lambda_-}=\sqrt{1-\mathcal I^2},
\]
\[
 \norm{\lambda_+-\lambda_-}^2=1-2C,\qquad
 \mathcal P=1-\sqrt{1-2C},\qquad
 \sum_j((\lambda_+)_j+i(\lambda_-)_j)^2=\mathcal I+2iC.
\]
Moreover, $0\le2C\le1$ and $0\le\mathcal P\le1$.
\end{lemma}
\begin{proof}
\leanok
\uses{def:balanced}
The two lobes partition a unit vector, so their squared norms sum to
one. Combine this with their difference $\mathcal I$, expand the
squared difference and the complex square, and use nonnegativity of
the lobe entries and of squared norms for the bounds.
\end{proof}

\begin{lemma}[Bimodal involution]
\label{lem:involution}
\lean{Appendices.iotaMap_iotaMap, Appendices.iotaMap_add, Appendices.iotaMap_smul, Appendices.sum_iotaMap_mul_iotaMap, Appendices.sum_iotaMap_mul, Appendices.memLambdaSet_iotaMap_iff, Appendices.lamPlus_iotaMap, Appendices.lamMinus_iotaMap, Appendices.imbalance_iotaMap, Appendices.statC_iotaMap, Appendices.pairingScore_iotaMap, Appendices.sum_sq_sub_iotaMap, Appendices.pairingScore_eq_dist_iotaMap, Appendices.pairingScore_eq_one_iff}
\leanok
\uses{def:balanced}
The map $\iota$ is an orthogonal, self-adjoint linear involution,
preserves $\Lambda_m$, and exchanges $\lambda_+$ and $\lambda_-$.
Consequently $\mathcal I\circ\iota=-\mathcal I$, $C\circ\iota=C$ and
$\mathcal P\circ\iota=\mathcal P$. For $\lambda\in\Lambda_m$,
\[
 \norm{\lambda-\iota\lambda}^2=2\norm{\lambda_+-\lambda_-}^2,\qquad
 \mathcal P(\lambda)=1-\frac1{\sqrt2}\norm{\lambda-\iota\lambda}.
\]
In particular, $\mathcal P(\lambda)=1$ if and only if $\iota\lambda=\lambda$.
\end{lemma}
\begin{proof}
\leanok
\uses{def:balanced}
Reversing the coordinates twice gives the identity; reversing and
negating preserves inner products and the ordering and sign pattern
defining $\Lambda_m$. Reindex the finite sums to obtain the
transformation rules. Split the displacement sum into its two halves
to identify both with the squared distance between the lobes.
\end{proof}

\setcounter{theorem}{3}
\begin{theorem}[Fibers of the profile map over the simplex faces]
\label{thm:faces}
\lean{Appendices.profileA_nonneg, Appendices.profileB_nonneg, Appendices.profileC_nonneg, Appendices.profileD_nonneg, Appendices.profile_sum_eq_one, Appendices.profileA_eq_zero_iff, Appendices.profileA_eq_one_iff, Appendices.profileB_eq_zero_iff, Appendices.profileB_eq_one_iff, Appendices.profileD_eq_zero_iff_forall, Appendices.profileD_eq_zero_iff_decomp, Appendices.profileD_eq_zero_decomp_commute, Appendices.profileD_eq_one_iff, Appendices.profileC_eq_zero_iff_forall, Appendices.profileC_eq_zero_iff_decomp, Appendices.profileC_eq_zero_decomp_commute, Appendices.profileC_eq_one_iff}
\leanok
\uses{def:profile}
Let $L\ne0$ and $\pi(L)=(a,b,c,d)$. The coordinates are nonnegative
and sum to one. Their faces and vertices have the following fibers.
\begin{enumerate}
\item $a=0$ exactly when $L$ is symmetric, and $a=1$ exactly when
$L$ is antisymmetric.
\item $b=0$ exactly when $S$ is semidefinite, and $b=1$ exactly when
$L$ is symmetric and $\alpha=\beta$.
\item $d=0$ exactly when $\alpha_j\ge\beta_j$ for every $j$,
equivalently $S=H+P$ with $H$ symmetric with spectrum symmetric about
zero and $P$ positive semidefinite. The matrices may be chosen to
commute with $S$ and with each other, with
$\norm{P}^2=c\norm{L}^2$. Also, $d=1$ exactly when $L$ is symmetric
negative semidefinite.
\item $c=0$ exactly when $\beta_j\ge\alpha_j$ for every $j$,
equivalently $S=H-P$ with the same conditions on $H,P$. They may be
chosen to commute with $S$ and with each other, with
$\norm{P}^2=d\norm{L}^2$. Also, $c=1$ exactly when $L$ is symmetric
positive semidefinite.
\end{enumerate}
\end{theorem}
\begin{proof}
\leanok
\uses{def:profile,lem:energy}
The spectral theorem and the energy partition give nonnegativity and
sum one. Vanishing squared norms characterize the $a$ face and the
coordinatewise inequalities for $c,d$. The lobe inner product vanishes
exactly when one lobe is zero, giving the $b$ face. When
$\alpha\ge\beta$, pair positive and negative eigenvectors to construct
$H$ with paired eigenvalues; the remainder $P$ is positive semidefinite.
Both are diagonal in an eigenbasis of $S$, giving commutation and the
norm formula. Conversely, adding a positive semidefinite matrix
increases the ordered eigenvalues, so comparison with the paired
spectrum of $H$ gives lobe dominance. Apply the argument to $-L$ for
the other face. The vertices follow from the partition of one.
\end{proof}

\chapter{Appendix C: Moments and the Gram form}

\begin{definition}[Odd, even and resolvent statistics]
\label{def:moment-statistics}
\lean{Appendices.statO, Appendices.statE, Appendices.statR}
\leanok
For a vector $\lambda$ with $|\lambda_i|<1$, define
\[
 O(\lambda)=\sum_i\frac{\lambda_i}{1-\lambda_i^2},\qquad
 E(\lambda)=\sum_i\frac{\lambda_i^2}{1-\lambda_i^2},\qquad
 R(\lambda)=\sum_i\frac1{1-\lambda_i}.
\]
\end{definition}

\begin{lemma}[Equivalent expressions for the moments]
\label{lem:moment-identities}
\lean{Appendices.MemLambdaSet.abs_entry_lt_one, Appendices.statO_eq_tsum_odd_moments, Appendices.statE_eq_tsum_even_moments, Appendices.summable_abs_odd_moments, Appendices.summable_abs_even_moments, Appendices.statO_eq_half_statR_sub, Appendices.statE_add_card_eq_half_statR_add, Appendices.statO_iotaMap, Appendices.statE_iotaMap, Appendices.trace_eq_statO, Appendices.trace_eq_statE, Appendices.trace_eq_statR}
\leanok
\uses{def:balanced,def:moment-statistics}
Every $\lambda\in\Lambda_m$ has $|\lambda_i|<1$. Writing $n=2m$,
\[
 O(\lambda)=\sum_{r=0}^{\infty}\sum_i\lambda_i^{2r+1},\qquad
 E(\lambda)=\sum_{r=0}^{\infty}\sum_i\lambda_i^{2r+2},
\]
and both series converge absolutely. Moreover,
\[
 O=\tfrac12(R-R\circ\iota),\qquad
 n+E=\tfrac12(R+R\circ\iota),\qquad
 O\circ\iota=-O,\qquad E\circ\iota=E.
\]
If a symmetric matrix $A$ has spectrum $\lambda$, with multiplicity, then
\[
 O(\lambda)=\tr(A(I-A^2)^{-1}),\quad
 E(\lambda)=\tr(A^2(I-A^2)^{-1}),\quad
 R(\lambda)=\tr((I-A)^{-1}).
\]
\end{lemma}
\begin{proof}
\leanok
\uses{def:balanced,def:moment-statistics,lem:involution}
The unit-norm condition and the presence of both signs force each
coordinate to have absolute value below one. Sum the convergent
geometric series coordinatewise. Partial fractions and reversal of
the coordinates give the identities involving $\iota$. For the trace
formulas, diagonalize $A$ orthogonally and evaluate the rational
functions on its eigenvalues.
\end{proof}

\begin{lemma}[Computation from the key and query matrices]
\label{lem:gram}
\lean{Appendices.GramForm.M, Appendices.GramForm.J, Appendices.GramForm.G, Appendices.GramForm.symPart_eq_half_conj, Appendices.GramForm.trace_pow_symPart, Appendices.GramForm.frobSq_symPart_eq_quarter_trace, Appendices.GramForm.charpoly_roots_filter_ne_zero, Appendices.GramForm.card_charpoly_roots_symPart, Appendices.GramForm.card_charpoly_roots_half_JG, Appendices.GramForm.IsSortedNormalizedNonzeroSpectrum, Appendices.GramForm.exists_isSortedNormalizedNonzeroSpectrum, Appendices.GramForm.isSortedNormalizedNonzeroSpectrum_unique, Appendices.GramForm.memLambdaSet_of_isSortedNormalizedNonzeroSpectrum, Appendices.GramForm.statO_eq_trace, Appendices.GramForm.statE_eq_trace, Appendices.GramForm.statR_eq_trace}
\leanok
\uses{def:parts,def:balanced,def:moment-statistics}
For $W_K,W_Q\in M_{n\times N}(\R)$, set
\[
 M=\begin{pmatrix}W_K\\W_Q\end{pmatrix},\qquad
 G=MM^T,\qquad J=\begin{pmatrix}0&I_n\\I_n&0\end{pmatrix}.
\]
The symmetric part of $W_K^TW_Q$ is $S=\tfrac12M^TJM$. For every
$k\ge1$,
\[
 \tr(S^k)=\tr((\tfrac12JG)^k),\qquad
 \norm{S}^2=\tfrac14\tr((JG)^2).
\]
If $M:\R^N\to\R^{2n}$ is surjective and $n\ge1$, the nonzero spectrum
of $S$, with multiplicity, is the spectrum of $\tfrac12JG$. Its
sorted normalized nonzero spectrum $\lambda$ exists uniquely and
belongs to $\Lambda_n$. For $H=JG/(2\norm{S})$,
\[
 O(\lambda)=\tr(H(I-H^2)^{-1}),\qquad
 E(\lambda)=\tr(H^2(I-H^2)^{-1}),\qquad
 R(\lambda)=\tr((I-H)^{-1}).
\]
\end{lemma}
\begin{proof}
\leanok
\uses{def:parts,def:balanced,def:moment-statistics,lem:moment-identities}
Block multiplication gives the factorization of $S$; cyclicity of
trace gives the power identities. Under surjectivity, $G$ is positive
definite. The characteristic-polynomial identity for rectangular
products identifies the nonzero spectrum, including multiplicities.
Congruence with $J$ gives exactly $n$ eigenvalues of each sign, so
sorting and normalization yield a vector in $\Lambda_n$. The matrix
$JG$ is similar to the symmetric matrix $G^{1/2}JG^{1/2}$; evaluate
the three rational functions on its real eigenvalues and take traces.
\end{proof}
